\chapter{Discussion and conclusions}
\label{ch:native-discussion}

\section{What the synthesis system delivers}
The first research question asks how to generate complete, class-conditioned journeys while respecting the available road topology and physical time. The implemented answer is a modular distribution over route, duration and progress. It accepts a map, OD request, participant class and public service constraints, and returns timestamped coordinates. The principal motion model is a trained diffusion denoiser. Learned route preferences make the delivered route distribution depend on observed choices rather than only on shortest paths or uniform sampling.

This factorization addresses the ambiguity of multiple valid paths between endpoints. A generated route is not judged incorrect merely because it differs from one SUMO realization. Route-use distributions, catalog coverage, movement distributions and within-condition diversity address distinct questions. The route experiment reduces edge-use divergence relative to shortest-route choice within its finite catalog; uniform route sampling retains advantages on other measures. The result therefore supports preference learning, not exact recovery of a unique future or exhaustive route generation.

The native clock addresses a different problem. Fixed-length relative-time profiles can lose physical resolution on long trips. The controlled temporal pilot substantially reduces native acceleration discrepancy when the diffusion target retains quarter-second timing, while longest-stop organization remains worse. The expanded native model then supplies a practical complete-journey baseline. The pilot is not a compute-matched large-corpus proof, and later results do not retroactively remove its limitations.

Road adherence and arrival follow from the chosen decoder. They are useful system properties, but should not be attributed to the GAF encoder or interpreted as complete traffic realism. The decoder has no separate lateral process, does not coordinate multiple generated agents and does not condition on a jointly evolving signal state. A plausible single trajectory can still be inconsistent with surrounding traffic or service dwell requirements. The output is suitable for studying synthetic-data distributions under the declared conditions, not a verified replacement for interactive traffic simulation.

\section{What GAF contributes, and what remains unproved}
The second research question asks whether participant-discriminative GAF representations improve synthesized trajectories. The inherited spatial method and the derived speed view both retain class information on native-demand data. The data-source diagnosis matters: randomly assigning classes and routes can remove behavioral associations that a spatial encoder is meant to recognize. Restoring native demand improves representation quality without adding velocity to the spatial input. This establishes a legitimate connection to the original classification work~\cite{salvadori,ferreira,original}, rather than redefining every image as a new GAF method.

Recognition is nevertheless only the first link. A deterministic image transform adds no observations beyond the coordinates and timestamps. It changes the inductive bias, compression and accessibility of existing information. Good classification can rely on route identity, service regularity or a broad speed scale, while generation needs the conditional distribution of motion after route, duration and class are already known. Information that is predictive before conditioning may be redundant after those conditions are supplied.

G40 makes this distinction testable. It adapts part of each encoder with the denoising objective, injects its code throughout sampling and models the code distribution using training data. This is a stronger connection than adding a classifier score only after generation, and follows the general encode--condition--prior idea of diffusion autoencoders~\cite{diffae}. The condition interventions show that the model uses the code. They do not show that the used information is sufficient, identifiable, or correctly sampled on new requests. The prospective physical results are the decisive check: the current dual-view implementation does not improve overall motion quality.

The spatial augmentation result is narrower and potentially useful. Synthetic J-S data help an independent classifier slightly more on average than B or RAW, under a fixed downstream budget. However, the effect changes direction across the two evaluation OD cohorts. It supports a cohort-dependent data-use observation, not a stable claim of better motion. The result should be replicated on new, prospectively fixed cohorts before becoming the central superiority claim of a final submission. A high recognition score or an attractive GAF image cannot substitute for that replication.

\section{Why the tested interface may remain limited}
Several explanations are consistent with the evidence, but none is isolated as the sole cause. First, the compact codes summarize entire journeys. Local stopping, acceleration and segment-specific behavior must pass through 32 values per view. The one-second speed construction and 128-position summary are deliberately insensitive to some fine motion changes. The sensitivity audit demonstrates this limitation directly; it does not prove that a higher-dimensional image code would solve it.

Second, reconstruction and prospective generation have different information. During joint training the encoder observes the target journey. At synthesis the code prior sees only the available route, duration and class. The small training-side gap between assigned and sampled J-SV codes is compatible with prior mismatch or target-specific information, while the short-prior pilot demonstrates a concrete instability. The longer formal prior removes the extreme range failure but does not establish correct conditional coverage. More latent updates alone should not be assumed to improve the delivered trajectories.

Third, the baseline already knows class, route geometry, length and duration. Adding another global class-related vector may largely repeat existing information. The class-centroid and same-class wrong-trip controls help expose this possibility. Their near-baseline or slightly favorable results prevent a strong causal story about correctly paired individual spatial/speed codes. The shuffle is not independent noise, and RAW differs in architecture and pretraining, so the present controls narrow the interpretation without uniquely identifying every mechanism.

Fourth, the joint objective combines noise prediction and class cross-entropy. A CE-only adaptation control is absent. Consequently the experiment cannot isolate whether differences between frozen and joint encoders arise from denoising gradients, classifier adaptation, or their interaction. Reported pre-adaptation classification figures should also not be presented as final adapted-encoder accuracy. These are explicit experimental limits, not reasons to reinterpret the adverse physical metrics as success.

\section{Evaluation must match the claim}
Physical quality is multidimensional. Travel distance alone misses stopping and acceleration; a native acceleration histogram misses ordering; a single reference's pointwise error can penalize another valid route. The thesis therefore retains class-marginal distances, request-level summaries, route statistics, coverage, diversity and an independent task. These measures sometimes disagree because they assess different properties. No post-hoc composite is introduced to turn a mixed outcome into an unqualified win.

Classifier confidence also requires caution. Calibration studies show why a predicted probability need not measure empirical correctness~\cite{calibration}, while proper scoring-rule theory distinguishes probability assessment from a task-specific reward~\cite{gneiting2007}. Even a calibrated participant classifier is not automatically a likelihood model for realistic trajectories within that class. A terminal score can favor a class stereotype or alter route/duration composition. The completed fixed-marginal controls and direct-generation comparison make those alternatives visible.

The experimental units must remain clear. Four OD cohorts with repeated departure windows are not eight independent domains. Two downstream seeds measure fitting variability on the same references; they do not independently replicate a geographic or behavioral population. The cohort breakdown of the spatial utility signal is therefore essential. Descriptive intervals, all-class tables and fixed display rules improve transparency, but cannot create additional independent evidence.

Observed references in this work are SUMO trajectories. The source alignment and native-demand expansion improve consistency with the supplied CNR material, yet they do not validate a transfer to field GPS traces. Simulator vehicle definitions, service schedules, map version and permission assumptions are part of the studied distribution. The second supplied CSV does not have a confirmed map match and is not silently mixed into Monaco evaluation. These boundaries protect the meaning of the resulting comparisons.

\section{A bounded next research step}
A justified next step would target information that remains unknown after conditioning, rather than repeat an unrestricted search for a favorable GAF weight. One candidate is a segment-level representation of route-conditioned motion, with spatial and speed codes trained to describe residual stop/go organization at an explicit physical clock. Another is joint modeling of travel time and progress, because the fixed cases expose duration mismatches that a fixed-duration motion intervention cannot correct. Both are proposals, not results of this thesis.

Any such extension should retain a matched raw-series encoder, frozen and class-only controls, a CE-only adaptation control, and a prior-versus-observed-code diagnostic on a separate training-side set. New evaluation cohorts must be fixed before viewing outcomes; shared bus services should be complemented by genuinely held-out services or routes if the claim is generalization. Downstream gain should recur across those cohorts and be read alongside physical quality and diversity. Increasing sample count without changing independent coverage would not resolve the current uncertainty.

\section{Conclusion}
This work delivers a reproducible road-conditioned diffusion pipeline for complete timed traffic journeys, with native-demand data preparation, learned route preferences, independent SUMO references and explicit information boundaries. It also implements and evaluates a substantial extension of inherited Hilbert--GAF classification: generation-trained spatial and speed conditions with a training-only latent prior.

The results answer the two research questions with different strength. The synthesis framework generates road-supported multi-class journeys and exposes meaningful temporal and routing tradeoffs. The GAF conditions contain recognizable class information and affect denoising, but the tested dual-view method does not improve overall motion realism. Spatial conditioning offers a small, cohort-dependent augmentation benefit, which remains distinct from physical fidelity. Keeping that distinction makes the contribution scientifically useful: it identifies the missing link between a representation that recognizes transport classes and one that reliably improves the conditional distribution of generated trajectories.
